
\subsection{Analytical strategy}

Log-linear models are the gold standard for analyzing marriage patterns, as they
allow us to estimate spousal associations while controlling for marginal
distributions of characteristics like ethnicity and education.
This control is particularly important because ethnic and educational
compositions of husbands and wives differ and change over time, and observed
marriage patterns are typically different from those expected under random
sorting.
Following convention, we refer to these spousal associations as assortative
mating preferences, though they reflect the causes to observed patterns
deviating from random ones rather than directly measured preferences.

We employ two sets of log-linear models.
The frist set estimates marriages by husband's and wife's ethncity (nine categories each)
across four census years (1980, 1990, 2000, and 2010), yielding 144 cells
($9 \times 9 \times 4 = 144$) after cross-tabulation.
The baseline model is specified as:

\begin{equation}
    \log F_{ijt} = \beta_0 + \beta_i^{HG} + \beta_j^{WG} + \beta_t^T + \beta_{it}^{HGT} + \beta_{jt}^{WGT}
\end{equation}

\noindent where $F_{ijt}$ is the number of marriages between husbands in ethnic
group $i$ and wives in ethnic group $j$ in year $t$,
$\beta_0$ is the constant,
$\beta_i^{HG}$ and $\beta_j^{WG}$ are husbands' and wives' ethnic groups,
respectively,
$\beta_t^T$ represents the year,
and $\beta_{it}^{HGT}$ and $\beta_{jt}^{WGT}$ denote the interactions between
year and ethnic group for husbands and wives, respectively.

In our second set of models, we incorporate educational attainment alongside
ethnicity.
These models predict the number of marriages for each combination of husband's
and wife's ethnicity and education across census years, yielding 2,916 cells
($9 \times 9 \times 3 \times 3 \times 4 = 2,916$).
The baseline model is specified as:

\begin{equation}
    \log F_{ijmnt} = \beta_0 + \beta_i^{HG} + \beta_j^{WG} + \beta_m^{HE} + \beta_n^{WE} +
    \beta_t^T + \beta_{imt}^{HGET} + \beta_{jnt}^{WGET}
\end{equation}

\noindent where $F_{ijmnt}$ is the number of marriages between husbands in
ethnic group $i$ and education group $m$ and wives in ethnic group $j$ and
education group $n$ in year $t$.
In addition to the previously defined parameters,
$\beta_m^{HE}$ and $\beta_n^{WE}$ denote husbands' and wives' education,
respectively,
and $\beta_{imt}^{HGET}$ and $\beta_{jnt}^{WGET}$ capture all two- and
three-way interactions between ethnic group, education, and year for husbands
and wives, respectively.

While log-linear models effectively disentangle assortative mating preferences
from marginal distributions, their use in analyzing status exchange within
minority-majority intermarriages has notable limitations:
they are complex, subject to controversial specifications and interpretations,
and unable to capture detailed status differentiation \cite{xie_new_2021}.
To address these limitations, we employ the Exchange Index,
a model-free approach proposed by \citeauthor{xie_new_2021}
\citeyear{xie_new_2021} that offers straightforward and unambiguous
interpretation.

The Exchange Index computation involves three essential steps for each
gender-specific ethnic pairing:

\begin{enumerate}
    \item Converting educational attainment into percentile ranks within gender
          and 11-year moving cohorts,
    \item Performing full exact matching between treated (interethnic marriage)
          and control (intra-ethnic marriage) cases, based on husbands'
          educational percentile and age, then wives',
    \item Estimating the difference in spousal educational percentiles between
          intra- and interethnic marriages using weighted samples matched by
          husbands' and wives' characteristics.
\end{enumerate}

The resulting indices ($EI_{H}$ and $EI_{W}$) indicate the average difference in
spousal educational status between inter- and intra-ethnic marriages, from both
husbands' and wives' perspectives.
Had the majority status been privileged would the indicaes of Han individuals be
significantly positive and that of minorities be significantly negative, and
vice versa if the minority status was considered more desirable.